\section{Background}
\label{sec:background}

Let \(A\) be a state-preparation circuit such that
\[
  A\ket{0}
  =
  \sqrt{1-a}\ket{\psi_0}\ket{0}
  +
  \sqrt{a}\ket{\psi_1}\ket{1},
\]
where the last qubit marks the event of interest and \(a\in[0,1]\) is the
amplitude to be estimated. For the amplification level \(k\), the ideal
Grover-amplified success probability is
\[
  p_k(a)
  =
  \sin^2\!\left((2k+1)\arcsin\sqrt{a}\right).
\]
Maximum-likelihood amplitude estimation (MLAE) estimates \(a\) by fitting the
observed counts across a schedule of amplification levels \(K\), avoiding the
full quantum phase estimation register used in canonical QAE.

Grover--Rudolph state preparation \cite{grover2002creating} is relevant for
numerical integration because efficiently integrable probability densities can
be encoded into quantum amplitudes through controlled rotations. The rigorous
version used as a reference point in this project formalizes the dyadic
probability tree, proves exact preparation of the target measurement law, and
gives perturbation and sampling design rules for finite-precision rotations
\cite{falco2026grproof}. In this project, the first hardware campaign uses the
small SISC calibration family \(g_0,g_1,g_2\), which is rich enough to test
constant, affine, and quadratic angle encodings while still allowing exhaustive
local auditing before hardware submission.

\subsection{Dyadic grid notation}
\label{subsec:dyadic-grid-notation}

For \(n\) grid qubits, set
\[
  B_n=\{0,1\}^n,\qquad N=2^n.
\]
The bit string \(b=(b_0,\ldots,b_{n-1})\in B_n\) labels the integer
\[
  i(b)=\sum_{j=0}^{n-1} b_j\,2^{n-1-j}\in\{0,\ldots,N-1\}.
\]
The midpoint dyadic grid is
\[
  D_n^{\mathrm{M}}
  =
  \left\{
    x_i^{\mathrm{M}}=\frac{i+1/2}{2^n}: i=0,\ldots,2^n-1
  \right\}.
\]
Unless a quadrature rule is explicitly stated, the shorter notation \(D_n\)
means \(D_n^{\mathrm{M}}\) throughout the Grover--Rudolph probability-law
experiments. For quadrature-rule QAE experiments we keep the rule visible and
write, for example,
\[
  x_i^{\mathrm{L}}=\frac{i}{2^n},\qquad
  x_i^{\mathrm{M}}=\frac{i+1/2}{2^n},\qquad
  x_i^{\mathrm{R}}=\frac{i+1}{2^n}.
\]
The Simpson value is not a new grid \(D_n\); in this report it denotes the
standard weighted combination of left, midpoint, and right evaluations used in
the corresponding quadrature script. This convention prevents a collision
between the support of a finite probability law and the quadrature rule used
to approximate a continuum integral.

\subsection{Angle-structure hierarchy for encoded functions}
\label{subsec:angle-structure-hierarchy}

The encoding-complexity paper \cite{chinesta2026encoding} introduces a hierarchy
of classes that is directly relevant to the present experiments. For an
\(n\)-qubit grid function \(g_n:B_n\to[0,1]\), usually obtained by sampling a
continuum function \(g\) as \(g_n(b)=g(x_{i(b)}^{\mathrm{M}})\), define the
angle map
\[
  \Theta_{g,n}(b)
  =
  2\arcsin\sqrt{g_n(b)}
  =
  2\arcsin\sqrt{g(x_{i(b)}^{\mathrm{M}})},
  \qquad b\in B_n .
\]
The sampled function belongs to \(\mathcal{G}^{(d)}_n\) when
\(\Theta_{g,n}\) has a
multilinear expansion in the bits \(b_0,\ldots,b_{n-1}\) of degree at most
\(d\):
\[
  \Theta_{g,n}(b)
  =
  \sum_{S\subseteq\{0,\ldots,n-1\}} \widehat{\Theta}_S
  \prod_{j\in S} b_j,
  \qquad
  \widehat{\Theta}_S=0\quad\text{for } |S|>d .
\]
This degree is not the ordinary polynomial degree of \(g(x)\). It is the degree
of the rotation-angle structure seen by the amplitude oracle. Consequently,
smoothness or a low-degree formula in \(x\) does not by itself imply a
low-complexity encoding. The calibration functions used at the start of this
project are the first three nontrivial levels of this hierarchy on the
two-grid-qubit midpoint grid:
\[
  g_0\in\mathcal{G}^{(0)}_2,\qquad
  g_1\in\mathcal{G}^{(1)}_2,\qquad
  g_2\in\mathcal{G}^{(2)}_2.
\]
In the software this classification is checked by applying the finite
difference/Walsh--Hadamard transform to the sampled angle map. The same
diagnostic is used below to label later test functions whenever they are
interpreted as amplitude-oracle functions.

\subsection{Angle structure and quantum integrability}
\label{subsec:angle-structure-integrability}

The angle hierarchy is not only a naming convention for benchmark functions; it
is a natural candidate for discussing when a finite integration problem is
pertinent for a quantum implementation. In an amplitude-estimation workflow, a
test function \(g_n:B_n\to[0,1]\) is not encoded through its values directly,
but through the rotation-angle map
\[
  \Theta_{g,n}(b)=2\arcsin\sqrt{g_n(b)}.
\]
If \(\Theta_{g,n}\in\mathcal{G}^{(d)}_n\) with small \(d\), then the angle
oracle has low multilinear interaction order in the grid bits. This does not
by itself prove a hardware advantage, but it is strong structural evidence
that the amplitude-encoding layer is compatible with short controlled-rotation
decompositions. Conversely, a function that is elementary as a formula in
\(x\) can still induce a full-degree angle map after the nonlinear
\(2\arcsin\sqrt{\cdot}\) transformation and restriction to \(D_n\), making it
much less suitable for amplified QAE on current hardware.

This report therefore separates two related notions of ``efficiently
integrable'':
\begin{itemize}[leftmargin=*]
  \item \emph{QAE-efficient integration}, where the quantity to be integrated
        is itself loaded as an amplitude function. Here the class
        \(\mathcal{G}^{(d)}_n\) is directly relevant because it controls the
        structure of the angle oracle.
  \item \emph{Grover--Rudolph distribution integration}, where a probability
        law \(p\) is first prepared and classical test functions \(f\) are
        integrated by postprocessing measured frequencies,
        \(I_p(f)=\sum_i p_i f(x_i)\). Here the relevant quantum cost is the
        probability-tree preparation of \(p\), while the angle class of \(f\)
        remains useful as a diagnostic for deciding whether the same quantity
        should later be attempted through amplified QAE.
\end{itemize}

The working hypothesis suggested by the present experiments is that low
angle-degree functions are the appropriate first candidates for amplified QAE,
whereas high angle-degree or localized quantities should first be studied
through direct Grover--Rudolph distribution sampling. The hypothesis is
deliberately conditional: actual efficiency also depends on the dyadic
probability law, grid size, compilation, two-qubit depth, mitigation strategy,
and the classical MC/QMC baselines used for comparison.

\subsection{Finite algebraic probability model}
\label{subsec:algebraic-probability-model}

All Grover--Rudolph cases in this report are finite probability models on a
dyadic grid. For \(n\) grid qubits we use the midpoint grid defined in
Section~\ref{subsec:dyadic-grid-notation},
\[
  D_n=D_n^{\mathrm{M}}
  =
  \left\{x_i=\frac{i+1/2}{2^n}: i=0,\ldots,2^n-1\right\}.
\]
Following the algebraic-probability viewpoint of
\cite{falco2026algebraicprobability}, the finite classical model is the
commutative algebra of grid functions
\(\mathcal{A}_n\simeq\mathbb{C}^N\). Its primitive events are the idempotents
\(e_i\), where \(e_i(x_j)=\delta_{ij}\). A probability vector
\(p=(p_0,\ldots,p_{N-1})\), with \(p_i\ge 0\) and \(\sum_i p_i=1\), defines the
state
\[
  \varphi_p(f)=\sum_{i=0}^{N-1}p_i f(x_i),
  \qquad f\in\mathcal{A}_n.
\]
Thus a numerical integral, payoff, event probability, or moment is treated as a
quantity-of-interest functional
\[
  I_p(f):=\varphi_p(f),
\]
where \(f\) is a test function in the finite algebra \(\mathcal{A}_n\). This is
the terminology used throughout the rest of the report: we estimate
functionals of a finite probability state, not Born-rule expectation values of
self-adjoint quantum operators.

The associated quantum system has Hilbert space
\(\mathcal{H}_n=(\mathbb{C}^2)^{\otimes n}\) and matrix algebra
\(\mathcal{B}(\mathcal{H}_n)\). A Grover--Rudolph state-preparation circuit
\(A_p\) implements the pure state
\[
  A_p\ket{0^n}
  =
  \sum_{i=0}^{N-1}\sqrt{p_i}\ket{i}.
\]
If \(\Pi_i=\ket{i}\!\bra{i}\) is the computational-basis projection, the quantum
state \(\omega_p(X)=\bra{0^n}A_p^\dagger X A_p\ket{0^n}\) satisfies
\[
  \omega_p(\Pi_i)=p_i.
\]
This identity is the bridge between the classical probability law on
\(\mathcal{A}_n\) and the noncommutative quantum experiment on
\(\mathcal{B}(\mathcal{H}_n)\). Direct distribution experiments estimate the
whole vector \(p\). Quantity-of-interest experiments estimate
\[
  I_p(h)=\sum_{i=0}^{N-1}p_i h(x_i)
\]
from measured frequencies \(\hat{p}_i\). Event-amplitude experiments choose an
event \(E\subset D_n\), for example the upper-half event
\[
  E_{\mathrm{upper}}=\{x_i\in D_n: x_i\ge 1/2\},
  \qquad
  a=p(E_{\mathrm{upper}})=\sum_{x_i\in E_{\mathrm{upper}}}p_i,
\]
and then use QAE or MLAE to estimate the scalar amplitude \(a\). For a Grover
amplification level \(k\), the ideal marked probability is
\[
  p_k(a)=
  \sin^2\!\left((2k+1)\arcsin\sqrt{a}\right).
\]

\subsection{Grover--Rudolph case names}
\label{subsec:gr-case-names}

The name \(\texttt{gr\_<profile>\_n<n>}\) has a precise meaning throughout the
software and this report. It denotes a Grover--Rudolph probability law on the
midpoint grid \(D_n\). A nonnegative profile \(F:D_n\to\mathbb{R}_{\ge 0}\) is
sampled on the grid, and the probability vector is
\[
  p_i=\frac{F(x_i)}{\sum_{j=0}^{2^n-1}F(x_j)},
  \qquad i=0,\ldots,2^n-1.
\]
The circuit then prepares \(\sum_i\sqrt{p_i}\ket{i}\). For example,
\(\texttt{gr\_sin2\_n3}\) means \(n=3\), \(N=8\), and
\(F(x)=\sin^2(\pi x)\). Similarly, \(\texttt{gr\_affine\_n4}\) means \(n=4\),
\(N=16\), and
\[
  F(x)=0.11+0.57x,
  \qquad
  p_i=
  \frac{0.11+0.57x_i}{\sum_{j=0}^{15}(0.11+0.57x_j)}.
\]
For this affine case, the upper-half event used in the first amplified
Grover--Rudolph test has
\[
  a
  =
  \sum_{i=8}^{15}p_i
  =
  0.6803797468\ldots.
\]

The structural complexity of the angle encoding is tracked through:
\begin{itemize}[leftmargin=*]
  \item multilinear support, the number of nonzero angle coefficients;
  \item maximum degree, the largest interaction order in the multilinear angle
        expansion;
  \item the minimum \(m_1\) stratum, which records how localized the smallest
        support term is on the grid;
  \item the weighted canonical length \(L_{\mathrm{can}}\), used as a compact
        proxy for angle-encoding complexity.
\end{itemize}

These structural metrics follow the angle-structure viewpoint developed for
quantum numerical integration \cite{chinesta2026encoding}. They complement
backend-dependent compilation metrics such as depth, two-qubit depth, and
two-qubit gate count. The Phase 1 results show that compilation choices can
dominate the hardware behavior even when the logical circuit and the saved
transpiled QASM are both ideal-correct.

Finally, the experiment is also informed by two broader theoretical references.
The intrinsic gate-embedding perspective of \cite{falco2026liegates} motivates
tracking compilation as a mathematical object rather than a purely software
detail. The algebraic probability viewpoint of
\cite{falco2026algebraicprobability} gives the finite-dimensional language in
which classical measurement laws, quantum states, and noncommutative operations
can be discussed within one probabilistic framework.
